\hspace{3mm}

\centerline{\bf \Huge Теория}
\centerline{\bf \Huge Дискретная непрерывность}

Начнем с простой задачи. 

\textbf{Задача 1}. В ряд выписано несколько целых чисел, причем любые два соседних либо равны, либо отличаются на единицу. Первое число ряда равно –10, а последнее равно 10. Докажите, что в этой последовательности есть число 0.

\textbf{Решение.} По существу, эта задача говорит о следующем: начав с целого отрицательного числа и сдвигаясь на один, мы не можем попасть в положительное число, «перескочив» при этом ноль. Это утверждение выглядит совсем очевидно, но все же аккуратно докажем его. Найдем в нашем ряду первое неотрицательное число (по условию неотрицательные числа в ряду присутствуют, а значит, можно взять первое из них). Может ли оно быть больше нуля? Нет, так как в этом случае предыдущее число обязано было бы быть неотрицательным, но мы взяли первое неотрицательное число в этом ряду. Значит, первое неотрицательное число равно нулю.

\textit{Вопрос.} В каком месте предыдущего доказательства использовалось условие, что первое число ряда равно –10?

Похожим образом можно доказать более общее утверждение.

\textbf{Дискретная теорема о промежуточном значении.} Пусть имеется последовательность целых чисел, в которой любые два соседних числа или равны, или отличаются на единицу. Пусть a – первый элемент последовательности, b – последний. Тогда для любого целого числа c, лежащего между a и b, найдется элемент последовательности, равный c.

\textbf{Дискретная теорема о промежуточном значении (переформулировка для любителей слова "функция").} Пусть имеется функция f, определенная на натуральных числах $k, k + 1, \dotsc, n$ и принимающая целые значения. Известно, что для любого $x = k , k + 1,\dotsc, n -1$ значения $f(x)$ и $f(x+1)$ либо равны, либо отличаются на 1. Пусть $A =f(k), B = f(n)$. Тогда для любого целого $C$, лежащего между $A$ и $B$, найдется такое число $x$, что $f(x) = C$.

Разберем еще несколько задач, применяя дискретную теорему о промежуточном значении.


\newpage

\hspace{3mm}

\textbf{Задача 2.} Первый тайм футбольного матча закончился со счетом 0:1, а второй тайм $-$ со счетом 4:3. Докажите, что в некоторый момент счет на табло был ничейный.

\textbf{Решение.} Рассмотрим последовательность, элементами которой являются разности между мячами, забитыми первой и второй командами. Первый элемент последовательности представляет собой разность в начале второго тайма, последний элемент – разность в конце второго тайма. Вместе с каждым забитым мячом появляется новый элемент последовательности (т.е. каждый элемент последовательности соответствует обновлению счета на табло). Так как первый элемент последовательности равен –1, последний равен 1 и каждые два соседних значения отличаются на 1, по дискретной теореме о промежуточном значении есть элемент последовательности, равный нулю. Этот элемент соответствует моменту, когда счет на табло был ничейным.

\textbf{Задача 3. Непростая!} В некоторой стране человек считается богатым, если его зарплата больше зарплаты премьер-министра. В этой стране богатые мужчины предпочитают жениться на бедных женщинах. Докажите, что можно установить премьерминистру такую зарплату, чтобы количество богатых мужчин было в точности равно количеству бедных женщин. (Все зарплаты в стране различные. Считайте, что зарплата бывает любым вещественным числом.)

\textbf{Решение.} Упорядочим по возрастанию все имеющиеся в стране зарплаты: $z_1, z_2, \dotsc, z_n$, зарплату премьер-министра обозначим через $Z$. Определим последовательность натуральных чисел так: $a_k$ равно разности между количеством богатых мужчин и бедных женщин, если $z_k < Z < z_{k+1}$. Кроме того, $a_0$ (соответственно, $a_n$ ) равно той же разности, если $Z < z_1$ (соответственно, $Z > z_n$). При таком определении любые два соседних элемента последовательности $a_k$ и $a_{k + 1}$ отличаются на $1$, так как при соответствующем изменении зарплаты премьер-министра человек с зарплатой $z_{k+1}$ становится богатым, что уменьшает текущую разность на $1$, если этот человек был женщиной, или увеличивает на 1, если он был мужчиной. Первый элемент последовательности соответствует ситуации, когда все люди в стране богатые (кроме самого премьера), и тогда $a_0 > 0$. Соответственно, последний элемент последовательности означает ситуацию, когда все люди бедные, и $a_n < 0$. По дискретной теореме о промежуточном значении существует такое k, что $a_k = 0$. Тогда установим зарплату премьер-министру, соответствующую этому значению 4.

\newpage

\hspace{3mm}

\textbf{Задача 4. Важная задача комбинаторной геометрии!} На окружности отмечены 77 точек, среди которых нет диаметрально противоположных. Докажите, что можно провести диаметр через одну из этих то- чек так, что по обе стороны диаметра точек окажется поровну.

\textbf{Решение.} Проведем диаметр окружности через произвольную отмеченную точку и раскрасим одну из половинок круга в синий цвет, а другую – в красный. Если в половинках поровну точек, то нужный диаметр найден. В противном случае определим последовательность следующим образом: $a_1$ равно разности между количеством отмеченных точек в синей части и красной части для данного положения диаметра. Будем постепенно поворачивать диаметр вокруг центра окружности (раскрашенные половинки круга поворачиваются вместе с ним) и каждый раз, когда с диаметра уходит отмеченная точка либо попадает на него, мы записываем новый элемент последовательности (который тоже равен разности между количеством отмеченных точек в синей части и красной части). Последний элемент последовательности соответствует такому положению диаметра, когда он снова проходит через первую отмеченную точку. В этом положении диаметр поворачивается на $180^{\circ}$ , а синяя и красная половинки меняются местами. Следовательно, последний элемент последовательности противоположен первому, т.е. один из них положителен, а второй – отрицателен. Так как любые два соседних элемента последовательности отличаются на 1, по дискретной теореме о промежуточном значении найдется элемент последовательности, равный нулю. Осталось заметить, что этот элемент последовательности соответствует положению диаметра, проходящему через одну из отмеченных точек, так как общее количество точек нечетно.